\section{Statement of the Problem}
\label{sec:problem}

\subsection{Background}

Modern deep learning workloads run close to the capacity limit of GPU memory.
The pressure comes from model parameters and activations, and also from
intermediate tensor transformations. Tensor layout conversion is a common
example.

Operations such as \texttt{permute}~\cite{pytorch_permute_docs} and
\texttt{transpose}~\cite{pytorch_transpose_docs} in
PyTorch~\cite{pytorch_2x} reorder the logical axes of a tensor without moving
any values. They only change the tensor's metadata, so they are free. What
they leave behind is a strided layout that is no longer contiguous.

Many downstream operations require a contiguous layout, and others merely run
slower on a strided one. A layout mismatch between what one operation produces
and what the next one needs is therefore routine. A recent PyTorch RFC reports
that non-contiguous tensors arise across sequence, tensor and pipeline
parallelism, and that several communication primitives require contiguous
inputs~\cite{pytorch_rfc_comms_2026}.

\subsection{The cost of the standard fix}

PyTorch resolves a mismatch with
\texttt{.contiguous()}~\cite{pytorch_contiguous_2016}. It allocates a new
tensor and copies every element into its new position. The copy is attractive
because it is highly parallel. Each output position can be computed
independently, so the GPU can use all of its threads.

The cost is memory. The original tensor and the new one must both exist until
the copy finishes. A single conversion can therefore approach twice the
footprint of the tensor being converted. On a large tensor that is enough to
cause an out-of-memory failure, and the short-lived allocation also fragments
the GPU heap~\cite{rajbhandari2020zero,pytorch_rfc_comms_2026}.

The workaround turns a layout problem into a peak-memory problem. Our goal is
to remove the second buffer instead, by converting the tensor in place.

\subsection{The problem, formally}

Let $A$ be a tensor of rank $r$ with shape $(d_0, \ldots, d_{r-1})$, stored
contiguously in row-major order, and let $\pi$ be a permutation of the $r$
axis positions. The \emph{in-place permutation problem} is to rewrite the
buffer holding $A$ so that it holds the tensor $A^{\pi}$ with shape
$(d_{\pi(0)}, \ldots, d_{\pi(r-1)})$, contiguous in the same row-major order,
using auxiliary space $o(N)$ where $N = \prod_i d_i$.

The constraint that makes this hard is that no value may be overwritten before
it has been read. The classical solution is cycle
following~\cite{gustavson2012cycleleader,dummit2003abstract}. A permutation
decomposes into disjoint cycles, and rotating the values along each cycle
moves every element exactly once with no second buffer. Cycle following is a
poor fit for a GPU, because cycle discovery is sequential and the resulting
accesses are irregular~\cite{catanzaro2014decomposition,gomezluna2016inplace}.

\subsection{Objectives}

\begin{enumerate}
\item Decompose an arbitrary $\pi$ into a sequence of operations, each of
  which is a single in-place matrix transpose. The cycle-following machinery
  that is well understood in two dimensions then applies unchanged at every
  step.
\item Choose the decomposition that moves the fewest bytes, rather than
  accepting the first one found, by searching the space of decompositions
  under an explicit cost model.
\item Keep auxiliary space independent of $N$. The only per-step state is the
  cycle metadata, which must stay at kilobyte scale for tensors of any size.
\item Beat the memory ceiling of \texttt{.contiguous()} on real shapes, and be
  no slower than it where both fit.
\end{enumerate}
